% 08_reproducibility.tex — Reproducibility and Open Infrastructure
% ============================================================

\section{Reproducibility and Open Infrastructure}
\label{sec:reproducibility}

Reproducibility in computational imaging is not merely a best practice
--- it is a prerequisite for the field to mature from artisanal
reconstruction to industrial-grade imaging.  The \pwm{} enforces
reproducibility at every layer through three mechanisms: immutable run
records, cryptographic integrity, and open governance.

% ============================================================
%  RunBundle
% ============================================================
\subsection{RunBundle: The Immutable Run Record}

Every imaging experiment executed through the \pwm{} produces a
\texttt{RunBundle}~v0.3.0 --- a self-contained, immutable archive that
records:
%
\begin{itemize}
    \item \textbf{Inputs:} raw measurements, metadata, and the
          \texttt{ExperimentSpec} that parameterized the run;
    \item \textbf{Operator state:} the nominal operator
          $\mathbf{H}_{\mathrm{nom}}$, detected mismatch parameters, and
          the corrected operator $\mathbf{H}_{\mathrm{corr}}$ (if
          correction was applied);
    \item \textbf{Triad diagnosis:} the full \texttt{TriadReport}
          including dominant gate, evidence scores, and confidence
          intervals;
    \item \textbf{Correction trajectory:} the sequence of correction
          steps, intermediate residuals, and convergence diagnostics;
    \item \textbf{Outputs:} reconstructed images and volumes with
          per-pixel uncertainty estimates;
    \item \textbf{Compute consumed:} wall time, peak memory, GPU
          utilization, and energy estimate.
\end{itemize}
%
Every artifact within the RunBundle is hashed with SHA-256.  The
top-level manifest includes a Merkle root so that any tampering with
individual files is detectable.  RunBundles are designed to be
self-describing: a future reader with access only to the bundle and the
\pwm{} codebase can reproduce the result without additional context.

% ============================================================
%  Decision Records
% ============================================================
\subsection{Decision Records for Imaging Systems}

We introduce \emph{Decision Records for Imaging Systems} (DR-IS),
inspired by architectural decision records in software engineering.
Every calibration decision --- whether to apply a correction, which
correction strategy to use, and what threshold triggered the decision ---
is cryptographically signed and logged within the RunBundle.  DR-IS
entries are append-only: corrections can be refined but never silently
overwritten.

% ============================================================
%  ExperimentSpec
% ============================================================
\subsection{ExperimentSpec: Eliminating Ambiguity}

The \texttt{ExperimentSpec}~v0.2.1, described in \cref{sec:agents}, is
the linchpin of reproducibility.  By requiring that every parameter be
drawn from a typed, validated registry, the \pwm{} eliminates the class
of irreproducibility caused by ambiguous or undocumented configuration.
When an LLM is used to generate experiment configurations, it returns
\emph{only} registry identifiers --- never raw numerical parameters ---
which are then resolved, validated, and frozen into the RunBundle.  No
free-form strings enter the pipeline.

% ============================================================
%  Open-core model
% ============================================================
\subsection{Open-Core Licensing and Contribution Model}

The \pwm{} evaluation harness is released under the MIT license,
ensuring that any researcher can audit, extend, and benchmark against
the framework without licensing barriers.  We define a four-level
contribution model:
%
\begin{enumerate}
    \item \textbf{Use the harness.}  Run the benchmark suite on your own
          data and methods.  No contribution required.
    \item \textbf{Submit methods.}  Register a new solver in the YAML
          registry and submit benchmark results via a pull request.
    \item \textbf{Contribute modalities.}  Implement a new forward
          operator, mismatch model, and correction strategy; submit with
          a TriadReport and RunBundle as evidence.
    \item \textbf{Become a steward.}  Take ownership of a modality
          vertical --- maintain the operator, curate datasets, and
          review community submissions.
\end{enumerate}

% ============================================================
%  Test infrastructure
% ============================================================
\subsection{Test Infrastructure}

The \pwm{} codebase is guarded by \textbf{3,743~tests} spanning unit
tests for individual operators, integration tests for the agent
pipeline, regression tests for known mismatch--correction pairs, and
end-to-end tests that verify RunBundle integrity.  The test suite
achieves \textbf{0~failures} on the current release branch.  Continuous
integration runs the full suite on every commit, and any benchmark
submission that introduces a test failure is automatically rejected.
